\section{Introduction}
\label{sec:intro}

Modular blockchains separate execution from services such as settlement and
data availability (DA), which independently operated networks provide. Consider
a chain whose validity predicate requires fresh evidence from such layers: a
block is valid only if it passes the BFT checks and carries both a Bitcoin
settlement receipt and a receipt from a DA network such as
Celestia~\cite{albassam2019lazyledger}, each attesting a height within a
tolerance of the current Bitcoin tip ($\tau_{\text{BTC}}$) or of the block ($\tau_{\text{DA}}$). Formally,
\[
\mathit{valid}(B) = \mathit{BFTvalid}(B) \wedge \mathit{BTCreceipt}(B)
\wedge \mathit{DAreceipt}(B).
\]
If Bitcoin settlement is congested or censored, or the DA network is
unavailable, no proposal satisfies the predicate. The BFT consensus
layer~\cite{buchman2016tendermint,buchman2018gossip} built on top stalls
although the validator set is fully connected and correct. The chain inherits
the liveness of the layers it depends on.

Dropping the requirement restores progress but suspends the guarantees that
motivated the modular design: without a settlement checkpoint the chain cannot
claim externally anchored settlement, and without DA confirmation it cannot
claim that its recent data is externally available. The problem is therefore
not only to detect an outage. The protocol must make all validators agree on
when an external guarantee may be suspended, restrict the operations whose
safety depended on it, and define the evidence required before the guarantee
is restored.

\textbf{The gap.} Existing work addresses parts of this problem under other
failure models. Protocols that trade finality for availability adapt to
validator participation or consensus-network conditions, not to external
infrastructure~\cite{neu2021ebbflow,buterin2020gasper}. Rollup mechanisms
provide alternative paths when sequencing or data publication fails, but the
decision rests with users, a sequencer, a committee, or a contract on the
parent chain~\cite{figueira2025escapehatch,arbitrum_anytrust,optimism_altda}.
Oracle designs and ABCI++ vote extensions let validators bring externally
observed values into replicated state~\cite{cosmos_adr064,slinky_oracle}, but
the value goes to the application. None of them covers a native chain whose
validity depends on independently operated external layers that different
validators observe differently. There, validators must agree on the observation
and on the change in safety policy that follows from it
(Section~\ref{sec:related}).

\textbf{Why the decision belongs in consensus.} The relaxation of receipt
requirements and the withdrawal circuit breaker change two things that every
validator must agree on: whether a block is valid, and whether a withdrawal
executes. Bitcoin, DA, and peripheral-network status are observed locally and
asynchronously, so two honest validators may see different conditions at the
same instant. Suppose each validator applied a local policy to its own view.
Validators would disagree on validity, so proposals would collect nil votes and
rounds would repeat, or they would execute withdrawals differently and produce
diverging \texttt{AppHash} values. A local, non-deterministic observation
cannot drive replicated state directly. It must first become a deterministic
input that every validator evaluates identically.

\textbf{Design approach.} Engram keeps CometBFT's proposer selection,
prevote/precommit protocol, quorum intersection, and deterministic finality
unchanged, and modifies the object and the validity policy that this pipeline
carries. At height $h{-}1$, each validator signs its sensor snapshot as a vote
extension; the sensors are read only there. The proposer of height $h$ includes
the signed snapshots of a quorum (more than two thirds of the voting power) in a \emph{health
certificate} $C^{h}$. Every validator verifies $C^{h}$, recomputes a
deterministic weighted-median/majority aggregate $\mathsf{Agg}(C^{h})$, and
accepts a proposal only if its target state satisfies
$s^{h} = \delta(s^{h-1}, \mathsf{Agg}(C^{h}))$, where $\delta$ is the
transition function of a four-state finite-state machine (FSM) with states
\textsc{Anchored}, \textsc{Suspicious}, \textsc{Sovereign}, and
\textsc{Recovering}. FSM state is written only from a certificate that
CometBFT has already agreed on (Section~\ref{subsec:health-certificate}).
``Hybrid Adaptive Consensus'' denotes this adaptive layer inside the
consensus pipeline, not a replacement for the agreement protocol.

The FSM makes the chain's security mode explicit. In \textsc{Anchored}, blocks
require the Bitcoin and DA receipts. Sustained degradation moves the chain
through \textsc{Suspicious} into \textsc{Sovereign}, where the quorum keeps
ordering blocks, the receipt requirement of each unhealthy layer is relaxed,
and cross-chain withdrawals are locked. Return to \textsc{Anchored} requires
recovery hysteresis, DA backfill of the \textsc{Sovereign} interval, and a
validator-verified zero-knowledge proof over that interval
(Sections~\ref{subsec:fsm} and~\ref{subsec:reanchoring}; the prototype's proof
gate is qualified in Section~\ref{subsec:limitations}).

\textbf{Scope.} Engram handles degradation of Bitcoin, Celestia, and the
peripheral P2P interfaces while the CometBFT quorum remains live. It does not
address loss of the core quorum: like any BFT protocol, it needs at least
a quorum (more than two thirds of the voting power) to communicate after global stabilization time.
``P2P degradation'' in this paper always means partial, peripheral degradation
(Section~\ref{subsec:threat-model}).

To our knowledge, no published design combines the following for a modular
chain whose validity depends on external settlement or DA layers:
(i)~a quorum-backed set of validator observations from which every validator
deterministically derives external-layer health; (ii)~deterministic adaptation
of both block-validity requirements and withdrawal admissibility from that
health while the native quorum keeps ordering blocks; and (iii)~a recovery
condition that requires a proof over the degraded interval and its DA backfill.
Engram targets this combination.

\textbf{Contributions.}
\begin{itemize}
    \item \textbf{Quorum-backed external-health agreement.} A height-bound
    certificate of signed observations from a quorum,
    carried in vote extensions, and a deterministic weighted-median/majority
    aggregation rule. When Byzantine validators hold less than one third of the voting power, every numeric aggregate stays within
    the range of the honest observations in a well-formed certificate,
    whichever observations a Byzantine proposer includes or omits
    (Section~\ref{subsec:health-certificate}).
    \item \textbf{Dependency-aware adaptive execution.} A four-state FSM on the
    certificate, with asymmetric hysteresis and exponential backoff on repeated
    failed recovery, that changes the Bitcoin and DA receipt requirements
    without modifying CometBFT's agreement core. A withdrawal circuit
    breaker~\cite{mendonca2020resiliency} excludes withdrawal-marked transactions while
    the chain runs without external corroboration
    (Sections~\ref{subsec:fsm}--\ref{subsec:transitions}; Withdrawal Admission
    Safety theorem in Section~\ref{subsec:analytical}).
    \item \textbf{Proof- and backfill-gated recovery.} A Noir/UltraHonk circuit
    that proves, for up to 256 \textsc{Sovereign}-interval headers at constant
    proving and verification cost, header-chain continuity, that every recorded
    state is \textsc{Sovereign} or \textsc{Recovering} with the lock flag set,
    the binding of the interval to the last anchored root and to the new root,
    and the quorum signatures of the two boundary blocks. It does not check
    transitions against the transition function or show that no withdrawal
    executed. Longer intervals chain several proofs
    (Appendix~\ref{app:aggregation}). The exit condition also requires that the
    interval's data was published to the DA layer. Validators verify each proof
    inside block execution (Section~\ref{subsec:reanchoring}).
    \item \textbf{Layered formal specification.} A four-layer TLA+
    specification. Bounded model checking of each layer's safety and liveness
    properties found zero violations at $N{\le}4$; the largest single run
    explored \resFSMStatesM{} distinct states. In a \resAblationCount-mechanism ablation study,
    removing each mechanism reproduces a distinct predicted counterexample. The
    cross-layer refinement obligation is specified but open
    (Section~\ref{sec:formal}).
    \item \textbf{Live evaluation.} A 7-validator Cosmos SDK testnet against
    live Bitcoin regtest and Celestia networks with network- and
    application-level fault injection, comparing Engram against a strict receipt-gated
    baseline and a vanilla CometBFT control
    (Sections~\ref{sec:setup}--\ref{sec:evaluation}).
\end{itemize}

\begin{table}[h]
\caption{Research Questions}
\label{tab:rq}
\centering
\small
\begin{tabular}{@{}p{0.07\linewidth}p{0.85\linewidth}@{}}
\toprule
\textbf{RQ} & \textbf{Question} \\
\midrule
RQ1 & \textbf{Safety.} Under Byzantine voting power below one third (two of the seven testbed validators), do honest
validators stay in agreement on the FSM state and the certified health, and is no
withdrawal-marked transaction committed while the FSM is locked? Safety here is consensus safety; asset-level safety is outside the prototype's scope (Section~\ref{sec:discussion}). \\
RQ2 & \textbf{Liveness.} When Bitcoin settlement or DA becomes unavailable,
does Engram keep committing blocks where a strict receipt-gated baseline stalls,
and what does it give up: how many withdrawals commit before its circuit
breaker locks, compared with a baseline that never locks? A degraded
peripheral P2P view does not affect a baseline gated on receipts, so it is
examined under RQ1. \\
RQ3 & \textbf{Recovery.} Once peripherals recover, do the DA backfill and the
recovery proof return the chain to \textsc{Anchored}, and how far does
hysteresis reduce oscillation, at what cost in detection delay? \\
RQ4 & \textbf{Cost.} What overhead do the extended consensus object,
certificate validation, and recovery proof add, and where does it come
from: proposal size, validation time, or additional consensus rounds? \\
\bottomrule
\end{tabular}
\end{table}

The remainder of this paper is organized as follows. Section~\ref{sec:related}
situates Engram in the literature. Section~\ref{sec:model} states the system
model, the threat model, and the security goals. Section~\ref{sec:protocol}
defines the protocol: sensors, the health certificate, the FSM, proposal
validity, and recovery. Section~\ref{sec:formal} states the security properties
and reports the formal specification and bounded model checking.
Section~\ref{sec:setup} describes the implementation, testbed, baselines, and
methodology, and Section~\ref{sec:evaluation} answers RQ1--RQ4.
Section~\ref{sec:discussion} discusses the security envelope, limitations, and
future work, and Section~\ref{sec:conclusion} concludes.
